%!TeX root=../tese.tex
%("dica" para o editor de texto: este arquivo é parte de um documento maior)
% para saber mais: https://tex.stackexchange.com/q/78101

\chapter{Metodologia}
\label{chap:metodologia}

\section{Delineamento da pesquisa}

Esta pesquisa caracteriza-se como um estudo empírico de engenharia de \textit{software}, com abordagem quantitativa e cuja finalidade é desenvolver e avaliar um método de vinculação de \textit{e-mails} de submissão de \textit{patches} aos respectivos \textit{commit}. O objeto de análise é a trajetória de alterações submetidas às listas de desenvolvimento do kernel Linux, desde o \textit{e-mail} que contém o \textit{patch} até sua eventual incorporação, sob a forma de \textit{commit}, à árvore principal do projeto. Para tanto, o estudo trata a associação entre esses artefatos como um problema de vinculação de registros, no qual registros provenientes de duas fontes distintas devem ser comparados para determinar se representam a mesma alteração de código.

O desenho metodológico compreende etapas articuladas de coleta das mensagens e dos \textit{commits}; transformação e normalização dos dados; seleção dos e-mails elegíveis; extração de uma amostra de \textit{commits}; blocagem; classificação automática dos pares; classificação manual inicial; classificação independente por dois revisores; mensuração da concordância inter avaliadores; treinamento da regressão logística; calibração probabilística; e avaliação temporal em dados não utilizados no ajuste. A sequência foi definida para reduzir progressivamente o espaço de comparação, preservar evidências de correspondência, produzir uma verdade de referência e estimar, para cada par, uma probabilidade calibrada de integração.

A unidade básica de análise é o par formado por um \textit{e-mail} elegível e um \textit{commit} candidato. Cada par reúne evidências textuais, temporais, autorais, estruturais e processuais. O desfecho supervisionado Y assume valor 1 quando o \textit{e-mail} corresponde ao \textit{commit} e 0 quando não corresponde. Em vez de limitar o resultado a uma decisão binária, o modelo produzirá P(\textit{Integração}), isto é, a probabilidade, entre 0 e 1, de o \textit{patch} descrito no \textit{e-mail} ter sido integrado naquele \textit{commit} da árvore principal. A decisão operacional e o encaminhamento para revisão manual serão derivados dessa probabilidade e dos limiares definidos após a calibração.

\section{Coleta dos dados}
Foram utilizadas duas fontes complementares. A primeira foi o conjunto \textit{LKML5Ws}, que reúne mensagens provenientes das listas de desenvolvimento do kernel Linux. A versão 1.0.0 foi publicada em 15 de julho de 2026 e contém quase 30 milhões de mensagens provenientes de aproximadamente 350 listas, abrangendo cerca de três décadas de desenvolvimento. O conjunto é organizado em arquivos colunares e disponibiliza tanto os metadados das mensagens quanto campos derivados relacionados ao assunto, corpo, código e encadeamento das conversações (\cite{passos2026lkml5ws}).

A coleta dos \textit{e-mails} não abrangeu indiscriminadamente toda a base. Foram selecionadas as mensagens a partir de 1º de outubro de 2015. Esse recorte temporal foi aplicado para manter compatibilidade com o período coberto pelos \textit{commits} e limitar o volume de dados processado. A leitura dos arquivos foi realizada de forma incremental, por partições e lotes, a fim de reduzir o consumo de memória. Além disso, foi mantida a organização original por listas, preservando a possibilidade de identificar a origem de cada mensagem.

A segunda fonte corresponde ao histórico do repositório principal do kernel Linux. A coleta de \textit{commits} partiu de uma cópia completa do repositório oficial, preservando o histórico disponível e as relações de ancestralidade entre as alterações (\cite{torvalds_linux_source}). Para assegurar a reprodutibilidade, a extração deve ser registrada a data da coleta e o identificador da revisão mais recente incluída. O universo destinado às comparações foi restringido aos \textit{commits} com data a partir de 1º de julho de 2015. O início desse período antecede em três meses o recorte dos \textit{e-mails}, evitando a exclusão de casos próximos ao limite temporal em decorrência de diferenças entre as datas de autoria, envio e integração.

Os dados dos \textit{commits} foram originalmente obtidos em formato \textit{JSON} por linha. Cada registro continha um objeto estruturado com autoria, datas, identificação do responsável pela integração, \textit{hash}, arquivos modificados, mensagem, relacionamentos de ancestralidade e outros metadados. O conteúdo também permitia recuperar os índices dos objetos Git anteriores e posteriores à alteração. 
A coleta foi planejada para que as duas fontes representassem etapas complementares do mesmo processo. O LKML5Ws registra a submissão e a revisão anteriores à integração, enquanto o histórico Git registra as alterações efetivamente incorporadas. Entretanto, a correspondência entre os conjuntos não é direta: o message-id identifica a mensagem nas listas, ao passo que o \textit{hash} identifica o \textit{commit}. Por esse motivo, os registros foram transformados em estruturas comparáveis antes da aplicação dos critérios de vinculação.


\section{Tratamento e transformação dos dados}
O tratamento dos \textit{e-mails} iniciou-se pela padronização dos tipos das colunas. Identificadores e campos textuais foram convertidos para cadeias de caracteres; datas foram convertidas para um tipo temporal comum; e campos com múltiplos valores, como etiquetas do assunto e trechos de código, foram representados como listas. Em seguida, foram removidos registros com message-id vazio ou duplicado. O primeiro registro de cada identificador foi mantido, preservando-se a ordem de ocorrência. Essa decisão considera que o message-id deve identificar unicamente uma mensagem no sistema de correio eletrônico.

Do conteúdo de código foram extraídas as linhas iniciadas pelo marcador \texttt{\textit{index}}, que, no formato de diff do Git, registram os identificadores dos \textit{blobs} anterior e posterior e, quando disponível, o modo do arquivo. As ocorrências foram armazenadas sem o prefixo textual \texttt{index}, mantendo a estrutura ``\texttt{índice\_anterior..indice\_posterior modo}''. Duplicidades dentro da mesma mensagem foram eliminadas, mas a ordem das ocorrências foi preservada. A extração dessas informações permite comparar tecnicamente a transformação representada pelo \textit{patch} com a transformação registrada no \textit{commit}. O procedimento de seleção e transformação dos e-mails encontra-se materializado no anexo I \texttt{EmailsElegiveis.py}.
 
Os registros elegíveis foram separados em duas estruturas relacionadas pelo message-id. A primeira reuniu os metadados necessários à blocagem e à comparação: message-id, assunto sem etiquetas (\texttt{untagged\_subject}), etiquetas do assunto (\texttt{subject\_tags}), data (date), data informada pelo cliente (\texttt{cliente\_date}) e índices Git. A segunda preservou o conteúdo de código. Essa separação evitou a repetição de campos extensos durante operações que dependiam apenas de metadados, reduzindo o consumo de memória e o custo de leitura.

Os dados dos commits também passaram por limpeza e normalização. Caracteres Unicode inválidos foram substituídos para assegurar a codificação em UTF-8. Entidades HTML foram convertidas, e etiquetas de hiperligação foram removidas sem eliminar endereços eletrônicos delimitados por sinais de menor e maior. Os campos textuais que não exigiam quebras de linha tiveram espaços, tabulações e caracteres de nova linha normalizados. Já o texto principal da mensagem do commit preservou as quebras entre linhas não vazias, devido à relevância da estrutura textual para análises posteriores. 

De cada registro foram extraídos o autor, a data de autoria, o integrador, a data de integração, o \textit{hash}, os arquivos modificados e a mensagem. A primeira linha da mensagem foi armazenada separadamente, pois normalmente corresponde ao assunto técnico da alteração e constitui atributo central para gerar pares candidatos. Também foram extraídas as ocorrências dos \textit{trailers}: \textit{Link}, \textit{Fixes} e \textit{Signed-off-by}. O campo Link é particularmente relevante, pois pode apontar explicitamente para a mensagem correspondente no arquivo na plataforma lore.kernel.org.

A mensagem principal foi delimitada antes da primeira ocorrência de \textit{Fixes}: ou, quando esse marcador não estava presente, antes de \textit{Signed-off-by}. Esse procedimento separou a descrição da alteração dos metadados processuais acrescentados ao final do texto. As datas de autoria e integração foram convertidas para UTC e tiveram a informação de fuso removida após a normalização, permitindo o cálculo uniforme de intervalos. Os dados processados foram relacionados a uma estrutura adicional que continha os índices Git por \textit{commit}. O arquivo anexo II - CommitsProcessamento.py evidencia essas etapas de extração, limpeza, padronização de datas e consolidação dos índices.

Por fim, os conjuntos foram gravados em formato \textit{Parquet} com compressão, estatísticas por grupo de linhas e particionamento quando necessário. Essa escolha favoreceu leituras seletivas e processamento em grande escala. Foram realizadas verificações de integridade quanto à existência e à ordem das colunas, aos tipos dos campos, à ausência de identificadores nulos ou vazios e à unicidade dos message-id's. Também foi verificado se as estruturas de metadados e de código continham exatamente o mesmo conjunto de mensagens.


\section{Critérios de elegibilidade dos \textit{e-mails}}
A população elegível foi definida para reunir mensagens que provavelmente representam submissões de \textit{patches}, excluindo respostas, encaminhamentos, discussões sem alteração de código e mensagens administrativas. Inicialmente, exigiu-se que a mensagem apresentasse uma etiqueta de \textit{patch} no assunto. Esse requisito foi operacionalizado pelo indicador \texttt{has\_patch\_subject\_tag}, que deveria assumir valor verdadeiro. O critério identifica assuntos que seguem convenções como [PATCH], incluindo variações relacionadas a versões ou séries.

A presença dessa etiqueta, entretanto, não é suficiente. Respostas geralmente preservam o assunto da mensagem original e, portanto, podem continuar contendo [PATCH]. Para reduzir esse tipo de falso positivo, foram excluídas as mensagens com marcadores de resposta, como ‘Re:’ ou ‘Response’. Também foram considerados os metadados de encadeamento, especialmente \textit{In-Reply-To}, porque eles permitem reconhecer que uma mensagem responde a outra, ainda que o marcador textual tenha sido alterado. O objetivo foi priorizar a mensagem que iniciou a submissão, e não as manifestações posteriores do \textit{thread}. 

De modo semelhante, foram excluídas mensagens com marcadores de encaminhamento, como ‘Fwd:’, ‘Fw:’ ou ‘Forwarded’. Essas mensagens reproduzem conteúdo enviado anteriormente e não correspondem necessariamente a uma submissão efetuada conforme o fluxo regular de desenvolvimento. A exclusão reduz duplicidades e evita que uma cópia encaminhada seja interpretada como origem do \textit{patch}. Assim, tanto o indicador de resposta quanto o de encaminhamento deveriam ser falsos ou nulos. Assim, para atender a esse requisito, tanto o indicador \texttt{has\_response\_subject\_tag} como o indicador \texttt{has\_forward\_subject\_tag} devem apresentar o valor falso. 

Foi igualmente exigido um message-id não nulo e não vazio. Esse identificador funciona como chave da mensagem e permite relacioná-la aos arquivos da lista, às \textit{threads} e às referências Link encontradas nos \textit{commits}. Mensagens sem identificador não oferecem uma chave confiável para rastreamento e, por isso, foram excluídas.

Além dos critérios baseados no assunto e nos metadados, cada mensagem deveria apresentar evidência concreta de alteração. Foram aceitas mensagens cujo campo de código estivesse preenchido ou cujo corpo contivesse uma linha iniciada por diff --git. Esse marcador introduz, no formato produzido pelo Git, o conjunto de alterações referente a um arquivo. Um \textit{patch} do Linux combina uma mensagem descritiva com as modificações expressas em formato de diff unificado, estruturado em arquivos, \textit{hunks} e linhas adicionadas ou removidas (\cite{hoang2021patchnet}).

A disjunção entre código não vazio e presença de diff --git foi adotada para tolerar diferenças de extração. Em alguns registros, o conteúdo pode ter sido identificado e armazenado no campo estruturado de código; em outros, o marcador permanece apenas no corpo bruto. Exigir simultaneamente as duas condições aumentaria o risco de excluir submissões válidas por deficiência de processamento, e não por ausência do \textit{patch}.

Em conjunto, esses critérios excluem respostas que apenas citam um \textit{patch}, mensagens encaminhadas, discussões sem código, comunicações administrativas, mensagens sem identificador e parte dos falsos positivos produzidos pela permanência de [PATCH] no assunto de respostas. Ainda assim, a população resultante deve ser interpretada como conjunto de mensagens provavelmente elegíveis, e não como verdade definitiva. Uma resposta pode conter uma nova versão válida, enquanto uma mensagem inicial pode reproduzir um diff sem constituir submissão formal. Consequentemente, a elegibilidade funciona como filtro de alta sensibilidade, cuja precisão será avaliada na revisão manual.

\section{Seleção da amostra de \textit{commits}}
Após o processamento do histórico, foi extraída uma amostra aleatória simples de 1.380 \textit{commits}. A seleção utilizou uma semente fixa, de modo que a mesma amostra possa ser reproduzida em execuções posteriores. A amostragem sobre os \textit{commits}, em vez de diretamente sobre os pares, assegura que o estudo parta de alterações reconhecidamente integradas à árvore principal e busque retrospectivamente as mensagens que podem tê-las originado.

Cada elemento da amostra preservou o \textit{hash}, as datas, o autor, a primeira linha da mensagem, a mensagem completa, os arquivos, os índices Git e os trailers extraídos. Essa estrutura foi utilizada para buscar possíveis mensagens correspondentes na população elegível. Quando um \textit{commit} apresentava referência Link, o identificador final do endereço era comparado ao message-id das mensagens. Nos demais casos, a busca inicial baseava-se principalmente na correspondência entre o assunto sem etiquetas e a primeira linha da mensagem do commit.

A amostra de 1.380 registros também serviu como base para construir a verdade de referência. Todavia, um mesmo \textit{commit} pode gerar mais de um par candidato quando existem diferentes versões do \textit{patch} com o mesmo assunto ou quando múltiplas mensagens compartilham atributos semelhantes. Por essa razão, a quantidade final de pares submetidos à revisão manual pode ser superior ao número de \textit{commits} amostrados.

\section{Blocagem (\textit{blocking}) e método de Classificação Automática}

A blocagem corresponde à etapa destinada a reduzir o produto cartesiano entre \textit{e-mails} e \textit{commits}. Sem essa restrição, cada uma das milhões de mensagens teria de ser comparada com todos os \textit{commits} da amostra, o que produziria um espaço de busca computacionalmente impraticável. Neste estudo, a blocagem considerou exclusivamente a igualdade, desconsiderando diferenças entre letras maiúsculas e minúsculas, entre o assunto normalizado do \textit{e-mail} de submissão, representado pelo campo \texttt{untagged\_subject}, e a primeira linha da mensagem do \textit{commit}. Essa correspondência foi utilizada apenas para gerar pares candidatos, uma vez que assuntos iguais não comprovam, isoladamente, que os registros representam a mesma alteração (\cite{ferreira2021incivility}). 

Para os \textit{e-mails} selecionados na etapa de blocagem, o conteúdo integral das mensagens foi recuperado no Lore, serviço de arquivamento das listas de desenvolvimento do kernel Linux, mantido pela \textit{Linux Kernel Organization}. Essa recuperação permitiu obter o remetente registrado no campo \textit{From} e autor no campo \textit{Signed-off-by}, informação necessária para comparar a autoria da mensagem com a autoria do \textit{commit} candidato. O Lore disponibiliza arquivos pesquisáveis das listas, preserva as mensagens por meio do sistema \textit{public-inbox} e permite sua localização a partir do message-id (\cite{linuxkernel_lore}).

Após essa complementação, os pares candidatos foram submetidos ao método de classificação automática. Essa etapa amplia o procedimento empregado por \cite{ferreira2021incivility}, que vinculou e-mails de revisão a commits mediante três heurísticas: correspondência entre o assunto do e-mail e o assunto da mensagem do commit; compatibilidade de autoria e ordem temporal; e alteração do mesmo arquivo, considerando uma proporção relevante de linhas modificadas. No estudo citado, a combinação dessas evidências alcançou precisão de 98,51\% e revocação de 90,08\% na identificação dos vínculos. No presente trabalho, a recuperação dos campos \textit{From} e \textit{Signed-off-by} possibilitou incorporar a compatibilidade de autoria à classificação automática, juntamente com as demais evidências textuais, temporais e técnicas.

No método proposto, o primeiro parâmetro da classificação verificou a existência de um vínculo explícito entre os artefatos. Quando o message-id do e-mail correspondia ao identificador presente no campo \textit{Link} do \textit{commit}, e o assunto normalizado era igual à primeira linha da mensagem, o par era classificado como ``Integrado''. Essa condição foi interpretada como evidência de elevada certeza, pois combina uma referência processual explícita à mensagem de submissão com a correspondência textual do resumo da alteração. O procedimento implementado no anexo III - EmailsCommits.py utiliza conjuntamente o vínculo explícito e a compatibilidade do assunto para reconhecer esses casos. 

Em seguida, foram tratadas as situações nas quais diferentes \textit{e-mails} apresentavam o mesmo assunto. Quando outro e-mail, distinto daquele em avaliação, possuía message-id correspondente ao campo \textit{Link} do \textit{commit} e compartilhava o mesmo assunto, o par corrente era classificado como ``Não integrado''. Essa regra permite diferenciar versões sucessivas de uma submissão que preservaram o mesmo título, evitando que o \textit{commit} seja atribuído indistintamente a todas as mensagens da série. A classificação, nesse caso, significa que o \textit{e-mail} avaliado não corresponde ao \textit{commit} em análise, ainda que possa estar relacionado a outra versão da alteração.

Na ausência de referência explícita, a classificação automática foi complementada por três parâmetros: compatibilidade da autoria, intervalo temporal e correspondência dos índices Git. A autoria foi considerada compatível quando o nome do remetente do e-mail, depois da remoção do endereço eletrônico, coincidia com o nome do autor do \textit{commit}. O parâmetro temporal foi satisfeito quando a diferença entre a data do \textit{commit} e a data do \textit{e-mail} se encontrava no intervalo de -1 a 180 dias. O limite inferior admite diferenças residuais decorrentes da normalização dos fusos horários, enquanto o limite superior contempla processos prolongados de revisão e integração. Como autoria e proximidade temporal podem ocorrer em pares não correspondentes, esses parâmetros não foram interpretados isoladamente como prova do vínculo. O arquivo no anexo III - EmailsCommits.py aplica esses critérios aos \textit{commits} candidatos que compartilham o assunto da mensagem.

O terceiro parâmetro complementar consistiu na comparação dos índices Git. Em um \textit{patch}, esses índices representam, por meio dos identificadores abreviados dos objetos blob, o conteúdo do arquivo antes e depois da alteração. A coincidência entre os índices do \textit{e-mail} e do \textit{commit} constitui evidência técnica forte de que ambos descrevem a mesma transformação. Entretanto, a divergência não exclui necessariamente a correspondência, pois alterações na base de código, operações de \textit{rebase}, resoluções de conflitos e modificações efetuadas durante a revisão podem produzir objetos Git distintos.

Para operacionalizar a comparação, os índices anteriores e posteriores foram normalizados para os primeiros 11 caracteres e analisados juntamente com o modo dos arquivos. O parâmetro foi considerado satisfeito quando os conjuntos de índices e modos extraídos do \textit{e-mail} estavam contidos nos conjuntos correspondentes do \textit{commit}. A relação de contenção permite que o \textit{commit} apresente transformações adicionais, desde que preserve aquelas identificadas no \textit{patch}. Índices compostos exclusivamente por zeros, presentes em operações de criação ou exclusão de arquivos, receberam tratamento específico.

A classificação final resultou da combinação das evidências. O par foi classificado como ``Integrado'' quando existia correspondência explícita entre message-id e Link ou, na ausência desse vínculo, quando autoria, intervalo temporal e índices Git eram simultaneamente compatíveis. Foi classificado como ``Não integrado'' quando outro e-mail de mesmo assunto apresentava vínculo explícito com o \textit{commit}, quando não existia \textit{commit} candidato com o mesmo assunto ou quando os três parâmetros complementares eram incompatíveis. 

Desse modo, a blocagem e a classificação automática desempenharam funções distintas. A blocagem utilizou apenas o assunto para gerar os pares candidatos, enquanto a classificação combinou vínculo explícito, autoria, temporalidade e índices Git para atribuir os pares às categorias ``Integrado'' ou ``Não integrado''. Essa separação metodológica evita interpretar a igualdade do assunto como evidência conclusiva e amplia as heurísticas de Ferreira et al. (2021) por meio da inclusão dos índices Git como representação técnica da transformação efetuada no código. 

\section{Classificação inicial por um revisor}

Para cada \textit{commit} da amostra, foram constituídos pares com os \textit{e-mails} selecionados na etapa de blocagem. Todas as informações produzidas durante a blocagem e a classificação automática foram disponibilizadas ao revisor, com o propósito de fornecer uma visão integrada das evidências associadas a cada possível vínculo. Para cada par, foram apresentados o message-id, o assunto, a data, a autoria, o corpo do e-mail e os índices Git do \textit{e-mail}. Em relação ao \textit{commit}, foram fornecidos o \textit{hash}, a primeira linha e o conteúdo integral da mensagem, as datas de autoria e integração, os arquivos modificados e o campo \textit{Link},

Além dos dados extraídos das fontes originais, foram apresentados os resultados das comparações efetuadas na etapa automática. Esses resultados incluíram indicadores de correspondência do assunto, do vínculo explícito entre message-id e \textit{Link}, da autoria, do intervalo temporal e dos índices Git, bem como a classificação automática preliminar do par nas categorias ``Integrado'' ou ``Não integrado''. Tais resultados foram utilizados como evidências auxiliares e não como substitutos do julgamento humano. Quando diferentes \textit{e-mails} ou \textit{commits} compartilhavam o mesmo assunto, todos os candidatos relacionados eram apresentados ao revisor, permitindo a comparação entre versões e a identificação do vínculo mais provável. 

Os casos que continham referência explícita no campo \textit{Link} foram mantidos na avaliação, embora apresentassem forte evidência de correspondência. Essa decisão teve como finalidade verificar a confiabilidade desse atributo e identificar possíveis inconsistências nos dados. Da mesma forma, foram preservados pares com concordância parcial, evidências contraditórias ou ausência de correspondência automática. Desse modo, evitou-se que a avaliação fosse composta exclusivamente por casos simples e assegurou-se a representação de diferentes níveis de complexidade, necessária ao treinamento, à calibração e à avaliação da regressão logística.

A avaliação manual inicial foi conduzida por um revisor com três objetivos: aperfeiçoar as definições operacionais das categorias, identificar evidências não contempladas pelo procedimento automático e elaborar um protocolo para a etapa posterior de classificação independente. Embora tivesse acesso à classificação automática, o revisor deveria fundamentar sua decisão no conjunto das evidências disponíveis. Consequentemente, sua avaliação poderia confirmar ou modificar o resultado atribuído pelo método automático. 

A categoria ``Integrado” foi atribuída quando as evidências indicavam que o \textit{e-mail} avaliado dera origem ao \textit{commit} ou representava uma versão substancialmente equivalente à alteração incorporada. Entre as evidências consideradas estavam a correspondência entre message-id e Link, a compatibilidade dos índices Git, a equivalência do conteúdo técnico, a correspondência da autoria, a proximidade temporal e a presença das mesmas alterações em arquivos coincidentes. A igualdade isolada do assunto não foi considerada suficiente para confirmar a integração, uma vez que esse atributo havia sido utilizado apenas na blocagem e poderia permanecer inalterado em diferentes versões de uma mesma submissão. 

A categoria ``Não integrado” foi aplicada quando havia evidências positivas de que o \textit{e-mail} e o \textit{commit} não representavam a mesma alteração. Essa situação ocorria, por exemplo, quando outro e-mail com o mesmo assunto apresentava vínculo explícito com o \textit{commit}, quando as modificações técnicas eram incompatíveis ou quando as evidências autorais, temporais e estruturais indicavam consistentemente artefatos distintos. Nesse contexto, a categoria significava que o \textit{e-mail} não correspondia ao \textit{commit} avaliado, mas não permitia concluir que a mensagem não tivesse sido integrada por meio de outro \textit{commit}. Do mesmo modo, a ausência de correspondência no procedimento automático não foi considerada, isoladamente, evidência suficiente para atribuir essa classificação. 

Por fim, as dúvidas, as divergências em relação à classificação automática e as evidências não previstas identificadas durante a avaliação foram registradas de forma sistemática. Esses registros subsidiaram o aperfeiçoamento das definições das categorias ``Integrado'' ou ``Não integrado'', bem como a elaboração do manual destinado à posterior classificação independente por dois revisores. Dessa forma, a avaliação inicial desempenhou três funções complementares: verificar os resultados da classificação automática, aprimorar o protocolo de classificação e preparar uma etapa subsequente de análise mais consistente e reprodutível. 

\section{Classificação independente e concordância inter avaliadores}

Após a etapa de classificação inicial, os pares foram classificados independentemente por dois revisores. Ambos empregaram o mesmo protocolo, contendo as definições das categorias, a ordem recomendada para examinar as evidências e exemplos limítrofes. Nenhum revisor teve acesso à decisão do outro durante a classificação, evitando influência recíproca.

A concordância foi mensurada pelo coeficiente \textit{kappa} de Cohen, apropriado para dois avaliadores que classificam unidades em categorias nominais mutuamente exclusivas. Diferentemente da concordância percentual bruta, o \textit{kappa} desconta a concordância esperada ao acaso (\cite{cohen1960coefficient}) (Cohen, 1960). O coeficiente foi calculado pela diferença entre a proporção observada de concordância e a proporção de concordância esperada ao acaso, dividida pela diferença entre 1 e a proporção de concordância esperada ao acaso. Esta última foi estimada com base nas distribuições marginais das classificações atribuídas pelos dois avaliadores.

Além do valor global de \textit{kappa}, serão apresentadas a matriz de confusão entre avaliadores, a concordância percentual e a distribuição das categorias. Esses resultados são importantes porque a prevalência elevada de uma categoria pode afetar a interpretação do coeficiente. As discordâncias serão analisadas qualitativamente e posteriormente resolvidas por consenso. Se a concordância não for alcançada, será encaminhado a um terceiro avaliador. O rótulo consensual constituirá a verdade de referência final.



\section{Regressão logística e probabilidade calibrada de integração}

A etapa de modelagem utilizará regressão logística supervisionada como modelo principal para estimar a probabilidade de integração do \textit{patch} representado pelo \textit{e-mail} no \textit{commit} candidato. A escolha decorre de sua interpretabilidade: cada coeficiente permite examinar a direção e a intensidade da associação entre uma evidência e a chance de correspondência, mantendo o resultado em uma escala probabilística compreensível. O objetivo não será apenas classificar pares como integrado ou não integrado, mas produzir P(Integração) entre 0 e 1 para apoiar decisões automáticas, priorização de revisão manual e análise de incerteza.

A variável dependente será definida como Yij=1 quando o \textit{e-mail} i e o \textit{commit} j representarem a mesma alteração, segundo a verdade de referência consensual, e Yij=0 quando não representarem. O conjunto de preditores será composto por evidências geradas nas etapas anteriores. A especificação após ajustes foi:

logit[P(Yij=1)] = $\beta$0 + $\beta$1CoincideLink + $\beta$2CoincideAutor + $\beta$3Intervalo + $\beta$4CoincideIndices + $\beta$5Delta + $\beta$6SimilaridadeDiffs.

CoincideLink indicará a igualdade validada entre o message-id do e-mail e o identificador recuperado do trailer Link do commit. CoincideAutor registrará a compatibilidade entre remetente, Signed-off-by e autor do commit. Intervalo representará a diferença temporal entre a submissão e a integração, preferencialmente em dias e, se necessário, com transformação ou termos não lineares. CoincideIndices registrará a igualdade ou contenção dos índices Git. Delta medirá diferença temporal entre a data do commit e a data do e-mail de submissão de patch. SimilaridadeDiffs é uma medida complementar do conteúdo técnico, que indica se a similaridade cosseno entre Diff do commit e Diff do \textit{e-mail} é maior ou igual 90.

Variáveis binárias serão codificadas como 0 e 1; medidas contínuas serão mantidas em escala compatível e padronizadas quando isso facilitar a interpretação numérica. A multicolinearidade será examinada por correlações e pelo fator de inflação da variância.
Os coeficientes serão usados para mostrar o efeito de cada variável sobre a chance de correspondência entre o \textit{e-mail} e o \textit{commit}, com intervalos de confiança, sem interpretar associações como relações causais. Um coeficiente positivo indicará aumento das chances estimadas de correspondência, mantidos os demais preditores constantes; um coeficiente negativo indicará redução. CoincideLink e CoincideIndices devem apresentar forte capacidade discriminativa, mas serão mantidos no modelo juntamente com evidências não determinísticas para permitir estimativas quando esses campos estiverem ausentes ou tiverem sido alterados durante a integração.

A separação dos dados preservará grupos relacionados ao mesmo \textit{commit}. Os pares serão distribuídos entre os conjuntos de treino, validação e teste por uma divisão aleatória agrupada pelo identificador do \textit{commit}, de modo que todas as mensagens candidatas associadas ao mesmo \textit{commit} permaneçam na mesma partição. Essa restrição evita que versões ou candidatos relacionados ao mesmo resultado de integração apareçam simultaneamente no treinamento e na avaliação, reduzindo o vazamento de informações e produzindo estimativas mais realistas do desempenho do modelo em \textit{commits} não observados durante o ajuste. O conjunto de validação será usado para ajustar o calibrador. O conjunto de teste permanecerá intocado até a avaliação final. Como a blocagem tende a produzir muito mais pares negativos do que positivos, serão relatadas a prevalência das classes e métricas adequadas ao desbalanceamento, incluindo precisão, revocação, medida F1, área sob a curva precisão-revocação e matriz de confusão.

\subsection{Dimensionamento e divisão da amostra}

A amostra de aproximadamente 1.400 pares oferece volume compatível com a especificação ajustada de seis preditores principais, mas a suficiência será verificada pelo número de observações na classe menos frequente e pelo total efetivo de parâmetros, e não apenas pelo tamanho global. Serão contabilizados separadamente os pares correspondentes e não correspondentes, incluindo graus de liberdade consumidos por transformações, indicadores de ausência e interações. Como verificação conservadora, buscar-se-á manter pelo menos 20 observações da classe minoritária por parâmetro estimado; se esse requisito não for alcançado, o modelo será simplificado. 

A divisão será aleatória e agrupada pelo \textit{hash} do \textit{commit} (CommitHash). Como referência operacional, aproximadamente 70\% dos pares serão destinados ao conjunto de treino, 15\% ao conjunto de validação e calibração e 15\% ao conjunto de teste. A atribuição será realizada no nível do \textit{commit}: primeiro, os identificadores distintos serão distribuídos entre as três partições; depois, todos os pares associados a cada identificador acompanharão a partição atribuída ao respectivo \textit{commit}. Em consequência da variação do número de candidatos por \textit{commit}, os quantitativos finais de pares poderão diferir ligeiramente das proporções pretendidas.

A divisão deverá utilizar uma semente fixa e um procedimento reprodutível de particionamento por grupos, por intermédio da função \textit{GroupShuffleSplit} do pacote scikit-learn (\cite{pedregosa2011scikit}). Após a divisão, são verificadas a ausência de \textit{CommitHash} compartilhado entre treino, validação e teste, a quantidade de \textit{commits distintos}, o número de pares e a distribuição de Y em cada partição. O conjunto de treino será usado para estimar os coeficientes; a validação será reservada ao ajuste do calibrador e à definição dos limiares; e o teste permanecerá intocado até a avaliação final.


\subsection{Calibração probabilística} 

Embora a regressão logística produza valores entre 0 e 1, essas saídas não devem ser presumidas como probabilidades calibradas. O desbalanceamento, a amostragem de candidatos, e a especificação incompleta podem fazer com que a confiança prevista difira da frequência real de acertos. Por isso, será aplicada uma etapa explícita de calibração usando apenas previsões fora do conjunto empregado no treinamento do modelo.

A técnica principal será o \textit{Platt scaling} (\cite{platt1999probabilistic}). Embora a técnica tenha sido originalmente apresentada para calibrar saídas de máquinas de vetores de suporte (SVM), ela pode ser aplicada como método de pós-processamento a outros classificadores, desde que o calibrador seja ajustado com previsões provenientes de dados não utilizados no treinamento do modelo-base. Na pesquisa, o método utiliza a pontuação produzida pelo modelo original e a ajusta por meio de uma segunda regressão logística. Esse ajuste procura fazer com que as probabilidades estimadas representem melhor a frequência real de correspondências. A escolha é compatível com a amostra total de aproximadamente 1.400 pares e com um subconjunto de validação estimado em cerca de 210 observações, pois o método estima poucos parâmetros, preserva a ordenação dos escores e tende a ser menos suscetível a sobre ajuste. 

No procedimento de calibração, ss probabilidades inicialmente estimadas pelo modelo-base serão ajustadas para evitar os valores extremos 0 e 1 e depois convertidas para a escala de logit. Sobre essa informação será ajustada uma regressão logística, que aprenderá a relação entre as probabilidades brutas e os resultados observados na variável dependente Y. O calibrador resultante produzirá probabilidades calibradas para os conjuntos de validação e de teste, sem utilizar os resultados do teste durante seu ajuste.

O desempenho probabilístico será avaliado nos conjuntos de validação e teste, considerando separadamente as probabilidades brutas do modelo e as probabilidades obtidas após o \textit{Platt scaling}. Serão calculados o \textit{Brier Score}, o \textit{log loss}, o \textit{Expected Calibration Error}, o intercepto e a inclinação de calibração. O \textit{Brier Score} também será calculado para um modelo de referência que atribui a todos os registros uma probabilidade constante correspondente à prevalência da classe positiva no conjunto de treino.

Para o cálculo do \textit{Expected Calibration Error}, as probabilidades serão distribuídas em dez faixas entre 0 e 1. Para cada faixa, serão apresentadas a quantidade de registros, a probabilidade média estimada e a frequência observada de correspondências. O erro será calculado como a média ponderada das diferenças absolutas entre a probabilidade média e a frequência observada em cada faixa. O intercepto e a inclinação de calibração serão estimados por uma regressão logística que utiliza o logit das probabilidades como variável explicativa.

A capacidade de discriminação será avaliada por precisão, revocação, medida F1 e área sob a curva precisão-revocação. As classificações binárias serão geradas com o limiar definido no valor de 0,50, e será construída uma matriz de confusão para cada combinação de conjunto e tipo de probabilidade. As avaliações serão realizadas para as probabilidades brutas e calibradas da validação e do teste.


\subsection{Categorias de confiança e decisão operacional} 
As probabilidades calibradas do conjunto de teste serão classificadas em cinco categorias descritivas: não correspondência provável para probabilidades inferiores a 0,50; confiança baixa para valores entre 0,50 e 0,80; confiança moderada entre 0,80 e 0,95; confiança alta entre 0,95 e 0,99; e confiança muito alta para probabilidades iguais ou superiores a 0,99. O código não cria uma categoria específica denominada “Confirmada”.

As probabilidades brutas e calibradas serão adicionadas aos conjuntos de validação e teste. No conjunto de treino, será armazenada somente a probabilidade bruta, utilizada para diagnóstico e não para o ajuste do calibrador. As informações serão reintegradas ao conjunto de dados do estudo, juntamente com a identificação da partição de origem. A categoria de confiança será atribuída apenas aos registros do conjunto de teste. No anexo IV – RegressaoCalibracao.py estão os procedimento de ajuste e calibração do modelo logístico.


\section{Reprodutibilidade e ameaças à validade}
A reprodutibilidade será assegurada pelo registro das versões dos conjuntos de dados, revisão do repositório, períodos temporais, semente da amostragem, critérios de elegibilidade, regras de blocagem e limiares empregados. As contagens de registros em cada etapa serão documentadas em um fluxo de seleção, permitindo identificar quantas mensagens foram excluídas por cada condição.

Entre as ameaças à validade, destaca-se a possibilidade de o conjunto \textit{LKML5Ws} conter mensagens ausentes ou extraídas incorretamente. A anonimização também exige recuperação complementar da autoria para parte da análise. A blocagem pode excluir vínculos verdadeiros, especialmente quando assunto, autoria e índices foram substancialmente alterados. A amostra de aproximadamente 1.400 pares é adequada ao desenvolvimento inicial do modelo, mas sua informação efetiva depende da prevalência de correspondências, do número de parâmetros e da concentração de vários candidatos nos mesmos \textit{commits}. Como a amostragem parte de \textit{commits} integrados, ela não representa, por si só, toda a população de \textit{e-mails} rejeitados; por isso, conclusões sobre não integração e generalização para períodos futuros serão formuladas com cautela.

Outra limitação decorre da dependência entre versões de uma mesma série e entre diferentes \textit{e-mails} candidatos ao mesmo \textit{commit}. Esses pares não constituem observações independentes. A divisão agrupada por \textit{hasj} reduzirá o vazamento ao impedir que candidatos do mesmo \textit{commit} sejam distribuídos entre treino, validação e teste. Entretanto, mensagens relacionadas à mesma série podem estar associadas a \textit{commits} distintos e, portanto, ainda introduzir dependência residual. Esse aspecto será documentado e considerado na interpretação dos intervalos de confiança e dos resultados de calibração.

Em síntese, a metodologia articula filtragem de alta escala, critérios determinísticos, julgamento humano, regressão logística supervisionada e calibração probabilística. A blocagem torna o problema computacionalmente viável; a revisão independente fornece uma verdade de referência; o \textit{kappa} avalia a reprodutibilidade dos julgamentos; a regressão logística combina as evidências de forma interpretável; e o \textit{Platt scaling} ajusta a saída para que P(Integração) se aproxime da frequência observada de correspondências. O resultado esperado é uma tabela de vínculos entre \textit{e-mails} e \textit{commits} acompanhada das evidências, da classificação final, da probabilidade calibrada de integração e da categoria de confiança.
